\section{Hedonism}

This content comes from chapter 1 of \textit{The Fundamentals of Ethics} which
starts on page 57.

\subsection{Happiness and Intrinsic Value}

\begin{readingnotes}
	\textbf{Instrumental goods} are things that are valuable because of the good
	things they bring about. \textbf{Intrinsically valuable} things are things whose
	value does not depend on anything else that is good; it is worth pursuing for
	its own sake.\\

	\textbf{Hedonism} defines happiness as the combination of pleasure and the
	absence of pain.\\

	There are two types of pleasure \textbf{physical pleasure} and \textbf{attitudinal pleasure}
	(enjoyment). Where the first is the result of physical stimuli. Hedonists believe
	happiness is attitudinal pleasure (the positive attitude of enjoyment).
\end{readingnotes}

\subsection{The Attractions of Hedonism}

\begin{readingnotes}
	The ancient Greek philosopher Epicurus was the first hedonist. He argued that
	the most pleasant condition is one of inner peace. This comes from two sources:
	moderation in all physical matters, and intellectual clarity about what is truly
	important.
\end{readingnotes}
